\chapter*{Wprowadzenie}
\addcontentsline{toc}{chapter}{Wprowadzenie}
Projekt składa się z czterech tomów: \emph{Rachunek prawdopodobieństwa}, \emph{Procesy stochastyczne}, \emph{Analiza stochastyczna} oraz \emph{Statystyka matematyczna}.
Pierwszy tom buduje podstawy teorii na przestrzeniach probabilistycznych:
od metod zliczania zdarzeń elementarnych do opisu rozkładów zmiennych losowych.
W obecnej części tomu wprowadzimy również funkcje tworzące i charakterystyczne,
rodzaje zbieżności zmiennych losowych, słabe i mocne prawa wielkich liczb
oraz centralne twierdzenie graniczne. Następnie zbudujemy
warunkową wartość oczekiwaną z twierdzenia Radona--Nikodyma,
wykażemy jej własności projekcyjne w $L^2$ oraz rozszerzymy
teorię granic o koncentrację, schematy trójkątne,
metodę delty i zbieżność dystrybuant empirycznych.
Dalsze rozdziały zawierają teorię dyskretnych łańcuchów Markowa,
filtracji, czasów zatrzymania, martyngałów, podmartyngałów i nadmartyngałów.
W dowodach szczegółowo rozróżniamy ograniczoność czasu,
całkowalność oraz jednostajną całkowalność.

Przez \emph{model klasyczny} rozumiemy w szczególności model
Laplace'a oparty na skończonej liczbie jednakowo prawdopodobnych
wyników. Nie jest on alternatywą dla miary probabilistycznej,
lecz jej bardzo szczególnym przypadkiem. Zanim przejdziemy do
modelowania danych i nieznanych parametrów, oddzielimy aksjomaty
rachunku prawdopodobieństwa od interpretacji konkretnych doświadczeń.

Przyjmujemy znajomość podstaw analizy matematycznej, całki Lebesgue'a,
pojęcia przestrzeni mierzalnej i twierdzenia Radona--Nikodyma.
Dowody skupiają się na konstrukcjach i istotnych przejściach logicznych;
w miejscach wymagających zaawansowanej teorii miary zaznaczamy, z jakiego
uprzednio znanego twierdzenia korzystamy.

W tej części pozostajemy przy aksjomatycznym i klasycznym rachunku
prawdopodobieństwa. Wzór Bayesa wystąpi jako tożsamość wynikająca
z definicji prawdopodobieństwa warunkowego, lecz interpretację
bayesowską parametrów, rozkłady a priori i a posteriori,
wiarygodność i wnioskowanie statystyczne odłożymy na później.

Podstawowe pojęcia będziemy wprowadzać według trzech pytań: jaki
problem prowadzi do danej definicji, co ta definicja rzeczywiście
mówi o losowym doświadczeniu i jak wynikają z niej twierdzenia.
Obok rachunków skończonych zestawimy modele o nieskończenie
wielu wynikach, w których intuicja oparta wyłącznie na liczeniu
może być myląca.

